\subsection{Admissibility of a heat law}\label{sec:admiss}
A heat law can be wrong in several ways, and the tests are not interchangeable. \emph{Conservation} is built into the balance laws and says nothing about the constitutive law. \emph{Thermodynamic admissibility} asks whether the law is compatible with the second law. For materials with internal state variables the Clausius--Duhem inequality restricts the response functions \citep{ColemanGurtin1967}. For heat conduction with memory it restricts the relaxation functions, not merely the signs of coefficients \citep{ChenNunziato1974}. \emph{Causality} asks that the response vanish before the excitation, which for the transfer function means analyticity for $\operatorname{Re}s>0$. \emph{Stability} asks that no mode grow in time. \emph{Finite speed} is a property of the high-frequency limit and stands independent of the others. A model can pass one test and fail another.

For the template of \cref{sec:foundations} in a linear, homogeneous medium, two necessary conditions follow from the transfer function $G(s)=\bigl[k(1+\tau_Ts)+k^*(1+\tau_vs)/s\bigr]/(1+\tau_qs)$ of the uncoupled thermal mode. Both are derived here and verified with the admissibility check of \cref{app:repro}.

\paragraph{Mean dissipation.} For a harmonic temperature gradient the time-averaged product $-q\,\theta_{,x}$ is proportional to $\operatorname{Re}G(i\omega)$, and
\begin{equation}\label{eq:reG}
\operatorname{Re}G(i\omega)=\frac{k+k^*(\tau_v-\tau_q)+k\,\tau_q\tau_T\,\omega^2}{1+\tau_q^2\omega^2}.
\end{equation}
A heat flux that on average removes entropy-producing gradients, with the entropy flux taken as $q/T_0$, needs $\operatorname{Re}G\ge0$ at every frequency. For $\tau_q,\tau_T\ge0$ this reduces to
\begin{equation}\label{eq:admiss}
k+k^*(\tau_v-\tau_q)\ \ge\ 0 .
\end{equation}
The Green--Naghdi type-II law ($k=0$, all lags zero) gives $\operatorname{Re}G=0$ at every frequency: it is lossless, which agrees with the statement of Green and Naghdi that this theory involves no dissipation of energy and permits harmonic waves without damping \citep{Green1993}.

\paragraph{Temporal stability.} For a real wave number $\kappa$, put $K=\kappa^2$. The modes $e^{st}$ of the uncoupled thermal equation solve
\begin{equation}\label{eq:cubic}
\rho c\,\tau_q s^3+(\rho c+Kk\tau_T)s^2+K(k+k^*\tau_v)s+Kk^*=0 .
\end{equation}
The Routh--Hurwitz conditions for all roots in the open left half-plane are equivalent to
\begin{equation}\label{eq:rh}
\rho c\,\bigl[k+k^*(\tau_v-\tau_q)\bigr]+Kk\tau_T\,(k+k^*\tau_v)>0 ,
\end{equation}
so that for $\tau_T\ge0$ stability at all wave numbers requires the same inequality \eqref{eq:admiss} as the dissipation test. The strict statement applies for $K>0$, $\tau_q>0$ and $k^*>0$. For $k^*=0$ the cubic has a factor $s$ that comes from the thermal-displacement variable ($\dot v=\theta$) and not from the temperature equation, so the test applies to the remaining quadratic factor; $K=0$ is the uniform, conserved mode ($s=0$); and for $\tau_q=0$ the degree drops and the quadratic test replaces the cubic one. Equation~\eqref{eq:rh} was confirmed against the numerically computed roots of \eqref{eq:cubic} for 20\,000 random parameter sets, with no disagreement. \Cref{tab:admiss} applies both tests to the template models.

\begin{table}[ht]\centering\small
\caption{Admissibility of the template models for the uncoupled thermal mode (derived here; $\tau_q,\tau_T,\tau_v\ge0$). ``Finite speed'' refers to the high-frequency limit of the uncoupled mode.}\label{tab:admiss}
\begin{tabularx}{\textwidth}{@{}p{2.3cm}p{3.0cm}p{3.6cm}X@{}}\toprule
Model & Parameter choice & Dissipation and stability & Comment\\\midrule
Fourier & all lags $=0$, $k^*=0$ & $\operatorname{Re}G=k>0$; stable & Infinite speed.\\
Lord--Shulman & $k^*=0$, $\tau_T=0$ & $\operatorname{Re}G=k/(1+\tau_q^2\omega^2)>0$; stable & Finite speed $\sqrt{k/(\rho c\,\tau_q)}$.\\
Dual-phase-lag & $k^*=0$ & $\operatorname{Re}G>0$ for $\tau_T\ge0$; stable & No finite front in the first-order law (diffusive at high frequency).\\
Green--Naghdi II & $k=0$, lags $=0$ & $\operatorname{Re}G=0$; marginal (purely oscillatory) & Lossless, finite speed.\\
Green--Naghdi III & lags $=0$ & $\operatorname{Re}G=k>0$; stable & Diffusive at high frequency, no finite front.\\
Moore--Gibson--Thompson & $\tau_T=\tau_v=0$ & stable iff $k\ge k^*\tau_q$ (marginal at equality) & Finite speed; not admissible if $k^*\tau_q>k$. The inequality is derived here from the uncoupled thermal mode, and Quintanilla \citep{Quintanilla2019} reports that stability of the model depends on the constitutive parameters.\\
Three-phase-lag & general & stable iff \eqref{eq:admiss} (for $\tau_T\ge0$) & Parameters must be chosen within this region.\\
\bottomrule\end{tabularx}
\end{table}

These tests are necessary, not sufficient. A mode-by-mode check says nothing about the coupled problem, about energy-flux postulates other than $q/T_0$, or about well-posedness of initial-boundary-value problems. The entropy flux in the Green--Naghdi theories is a separate quantity, with the heat flux equal to the temperature times that flux. For type II it follows from the free-energy potential rather than from a Fourier law. In linearized type I it agrees with $q/T_0$ to first order \citep{Green1993}. For a law with memory the second law constrains the relaxation functions themselves. The initial values of the energy and heat-flux relaxation functions are non-negative, and the initial slope of the heat-flux relaxation function is non-positive \citep{ChenNunziato1974}. The equilibrium conductivity is non-negative \citep{ChenNunziato1974}. A Kramers--Kronig check of the kind in \cref{sec:prelim} does not establish admissibility. Inequality \eqref{eq:admiss} is also a parameter restriction that fitted values must respect, and reported phase-lag values should be checked against it before they are used as evidence for a model.
